vLLMは2月4日のv0.15.1を経て、2月25日にv0.16.0を公開した。後者のrelease noteはbranch cutが2月8日だったことも明記している。本号ではv0.15.1を月初のcontextとし、より強いsystems signalを持つv0.16.0を中心に読む。branch cutという時間境界がある以上、2月後半に登場した全model changeがv0.16.0へ既に取り込まれていた、とrelease日だけから推論することはできない。\autocite{src-5041feb6f87f,src-67307025f5d9}


v0.16.0で目立つのは、WebSocket-based Realtime APIとstreaming audio / Voxtral realtime support、Responses APIの改善、tool-calling fixes、unified parallel drafting、XPU関連の大きな変更が同じreleaseへ並んだことだ。これはserving engineがtext completionを高速に返すだけでは足りず、音声stream、tool protocol、agent-style response surface、複数のdecoding/scheduling経路を同時に扱う必要が出てきたことを示す。model capabilityの拡張がruntime APIの複雑さへ直結している。\autocite{src-5041feb6f87f}


